\chapter*{Part Synthesis: Models and Representations}
\addcontentsline{toc}{chapter}{Part Synthesis: Models and Representations}

A neural network composes transformations. The chain rule connects the final loss to each parameter, and optimization uses those derivatives to seek a fitted rule. The objective, initialization, step size, and model constraints all affect which rule is found.

Probability models add another perspective: parameters and latent variables can explain how observations arise, and predictive distributions can describe uncertainty within those assumptions. The law of total variance separates two contributions,
\[
\operatorname{Var}(Y\mid x,D)=
E_{\theta\mid D}[\operatorname{Var}(Y\mid x,\theta)]
+\operatorname{Var}_{\theta\mid D}[E(Y\mid x,\theta)].
\]
This decomposition does not protect a misspecified model from shared errors; it states what its uncertainty means.

Representation learning chooses an objective that preserves selected information. Reconstruction, masked prediction, and contrastive learning impose different constraints. Generative objectives describe distributions over possible observations. Useful transfer is a further empirical question: does the learned representation preserve the distinctions needed by the new task?

\clearpage
\begingroup\pagestyle{empty}\cleardoublepage\endgroup
